\begin{helper}
Let $\Delta$ be a positive integer and let $\gamma$ be a real number with
$0\le \gamma\le \Delta$.  Then
\[
\left|
\frac1\Delta\int_0^\gamma \left(1-\frac{x}{\Delta}\right)^\Delta\,dx
-\frac{1-e^{-\gamma}}{\Delta}
\right|
\le \frac2{\Delta^2}.
\]
\end{helper}
\begin{proof}
For every $x\in[0,\gamma]$ we also have $0\le x\le\Delta$.  The elementary
inequality $(1-t/n)^n\le e^{-t}$, applied with $n=\Delta$ and $t=x$, gives
\[
\left(1-\frac{x}{\Delta}\right)^\Delta\le e^{-x}
\]
on this interval.  Therefore the integral of the polynomial factor is at most
the integral of $e^{-x}$ over $[0,\gamma]$, and the difference
\[
\int_0^\gamma e^{-x}\,dx-
\int_0^\gamma \left(1-\frac{x}{\Delta}\right)^\Delta\,dx
\]
is nonnegative.

The same pointwise comparison is valid on the whole interval $[0,\Delta]$.
Since the difference of the two integrands is nonnegative there, and since
$[0,\gamma]\subseteq[0,\Delta]$, the preceding difference is at most
\[
\int_0^\Delta e^{-x}\,dx-
\int_0^\Delta \left(1-\frac{x}{\Delta}\right)^\Delta\,dx .
\]
Both endpoint integrals are explicit.  First,
\[
\int_0^\Delta e^{-x}\,dx=1-e^{-\Delta}.
\]
Second, with the substitution $u=1-x/\Delta$,
\[
\int_0^\Delta \left(1-\frac{x}{\Delta}\right)^\Delta\,dx
=\Delta\int_0^1 u^\Delta\,du
=\frac{\Delta}{\Delta+1}.
\]
Thus
\[
\int_0^\Delta e^{-x}\,dx-
\int_0^\Delta \left(1-\frac{x}{\Delta}\right)^\Delta\,dx
=(1-e^{-\Delta})-\frac{\Delta}{\Delta+1}
\le 1-\frac{\Delta}{\Delta+1}
=\frac1{\Delta+1}
\le\frac1\Delta .
\]
Since $\int_0^\gamma e^{-x}\,dx=1-e^{-\gamma}$, we have
\[
0\le
(1-e^{-\gamma})-
\int_0^\gamma\left(1-\frac{x}{\Delta}\right)^\Delta\,dx
\le \frac1\Delta .
\]
Multiplying by the positive factor $1/\Delta$ gives
\[
\left|
\frac1\Delta\int_0^\gamma \left(1-\frac{x}{\Delta}\right)^\Delta\,dx
-\frac{1-e^{-\gamma}}{\Delta}
\right|
\le \frac1{\Delta^2}\le \frac2{\Delta^2},
\]
which is the desired estimate.
\end{proof}
